% !TeX document-id = {b657b89f-3a7b-42a1-879d-53584fcd9464}
% !TeX spellcheck = en_GB
% !TeX TS-program = lualatex

\documentclass[final,aspectratio=169,12pt]{beamer}

\usepackage[english]{babel}

%\usetheme{plain}
%\usepackage[math-style=ISO,warnings-off={mathtools-colon,mathtools-overbracket}]{unicode-math}
%\usepackage{lualatex-math}
\usepackage[no-math]{fontspec}
\usepackage[defaultmathsizes, basic]{mathastext}
\usepackage{fontawesome5}
%\setmathfont{Latin Modern Math}
%\setmathfont{Source Sans Pro}
\setmonofont{Inconsolata SemiBold}
%\setmathfont[range=\symup]{Latin Modern Roman}
\DeclareMathAlphabet{\Lmathit}{\encodingdefault}{\familydefault}{m}{it}

\setmainfont{SourceSansPro}

\newfontfamily{\titlefont}{League Gothic Regular}[]
\usepackage{amsmath}
\usepackage{xcolor}
\usepackage{qrcode}
\usepackage{tikz}
\usetikzlibrary{matrix,positioning,calc,fit,shapes,scopes,arrows.meta,backgrounds,decorations.markings,3d,decorations.pathreplacing,node-families,overlay-beamer-styles}

% also works for colorblind
\definecolor{SeqGrey1}{HTML}{595959}
\definecolor{SeqGrey2}{HTML}{808080}
\definecolor{SeqGrey3}{HTML}{B1B1B1}
\definecolor{SeqGrey4}{HTML}{ECECEC}

\definecolor{SeqGreen1}{HTML}{003213}
\definecolor{SeqGreen2}{HTML}{277D35}
\definecolor{SeqGreen3}{HTML}{6AB969}
\definecolor{SeqGreen4}{HTML}{B6E6AE}

\definecolor{SeqBlue1}{HTML}{124066}
\definecolor{SeqBlue2}{HTML}{2279BF}
\definecolor{SeqBlue3}{HTML}{2EA1FF}
\definecolor{SeqBlue4}{HTML}{A7C7FD}

\definecolor{SeqYellow1}{HTML}{735805}
\definecolor{SeqYellow2}{HTML}{d1a009}
\definecolor{SeqYellow3}{HTML}{F7CB45}
\definecolor{SeqYellow4}{HTML}{fdf0c9}

\definecolor{SeqRed1}{HTML}{7c0307}
\definecolor{SeqRed2}{HTML}{DC050C}
\definecolor{SeqRed3}{HTML}{fb5b61}
\definecolor{SeqRed4}{HTML}{fdbbbd}

\definecolor{BuRdBlue1}{HTML}{2166AC}
\definecolor{BuRdBlue2}{HTML}{4393C3}
\definecolor{BuRdBlue3}{HTML}{92C5DE}
\definecolor{BuRdBlue4}{HTML}{D1E5F0}
\definecolor{BuRdMidpoint}{HTML}{F7F7F7}
\definecolor{BuRdRed4}{HTML}{FDDBC7}
\definecolor{BuRdRed3}{HTML}{F4A582}
\definecolor{BuRdRed2}{HTML}{D6604D}
\definecolor{BuRdRed1}{HTML}{B2182B}
\definecolor{BuRdBadData}{HTML}{FFEE99}

\definecolor{BgColor}{HTML}{f5f5f5}

\setbeamerfont{frametitle}{size=\Huge,family=\titlefont}
\setbeamerfont{title}{size=\Huge,family=\titlefont}
\beamertemplatenavigationsymbolsempty
%\setbeamertemplate{footline}[page number]
%\addtobeamertemplate{navigation symbols}{}{%
%    \usebeamerfont{footline}%
%    \usebeamercolor[fg]{footline}%
%    %\hspace{1em}%
%    \normalsize\insertframenumber
%}

% access relative slide number (animations)
\makeatletter
\def\c@slideinframe{\beamer@slideinframe}
\def\beamerslideinframe{\beamer@slideinframe}
\makeatother

\setbeamertemplate{footline}{
  \hfill%
  \usebeamercolor[fg]{page number in head/foot}%
  \usebeamerfont{page number in head/foot}%
  \setbeamertemplate{page number in head/foot}[framenumber]%
  \usebeamertemplate*{page number in head/foot}\kern1em\vskip2pt%
}
\setbeamersize{text margin left=10mm,text margin right=10mm}

\setbeamercolor{title}{fg=SeqGrey1}
\setbeamercolor{frametitle}{fg=SeqGrey1}
\setbeamercolor{normal text}{fg=SeqGrey1}


%\setbeamertemplate{frametitle}{\normalseries}[default][center]

\title{LEVERAGING FSPMs FOR UNCONVENTIONAL COMPUTING WITH PLANTS}
\date{30 March 2023}
\author{Olivier Pieters}
\institute{FSPM 2023}

\definecolor{input green}{HTML}{70dc70}
\definecolor{reservoir red}{HTML}{ff7d4d}
\definecolor{output blue}{HTML}{51a9ff}
\definecolor{weight olive}{HTML}{898900}
\definecolor{training purple}{HTML}{d5bfea}
\definecolor{dark purple}{HTML}{7339ac}
\definecolor{output pink}{HTML}{ef7bcc}
\definecolor{output blue}{HTML}{51a9ff}
\definecolor{light grey}{HTML}{afafaf}

\tikzstyle{<line>} = [{Latex[length=2mm, width=2mm]}-{Latex[length=2mm, width=2mm]},ultra thick]
\tikzstyle{line>} = [-{Latex[length=2mm, width=2mm]},ultra thick]
\tikzstyle{<line} = [{Latex[length=2mm, width=2mm]}-,ultra thick]
\tikzstyle{line} = [ultra thick]
\tikzstyle{input node} = [draw,circle,minimum width=0.75cm,line width=0.3mm,fill=input green]
\tikzstyle{reservoir node} = [draw,circle,minimum width=0.5cm,line width=0.3mm,fill=reservoir red]
\tikzstyle{output node} = [draw,circle,minimum width=0.75cm,line width=0.3mm,fill=output blue]
\tikzstyle{box} = [dashed,line width=0.3mm]
\tikzstyle{training weights} = [weight olive]
\tikzstyle{training} = [training purple,fill=training purple]
\tikzstyle{box label} = [font=\Large]

\tikzset{
	block/.style = {draw,
		rectangle,
		minimum height=3em,
		minimum width=3em},
	tmp/.style = {coordinate},
	sum/.style = {draw, circle, node distance=1cm},
	input/.style = {coordinate},
	design/.style = {fill=green,fill opacity=0.1,text opacity=1},
   overlap label/.style={fill=white,fill opacity=0.85,text opacity=1,text=fg,},
	output/.style = {coordinate},
	lowpass/.pic={%
		\draw (-0.4cm,0.25cm)
		-- ++(0.5cm,0)
		-- ++(0.3cm,-0.5cm);
		\draw (-0.4cm-0.1cm,0.1cm+0.4cm)
		-- (0.4cm+0.1cm,0.1cm+0.4cm)
		-- (0.4cm+0.1cm,-0.1cm-0.4cm)
		-- (-0.4cm-0.1cm,-0.1cm-0.4cm)
		-- cycle;},
	highpass/.pic={%
		\draw (0.4cm,0.25cm)
		-- ++(-0.5cm,0)
		-- ++(-0.3cm,-0.5cm);
		\draw (-0.4cm-0.1cm,0.1cm+0.4cm)
		-- (0.4cm+0.1cm,0.1cm+0.4cm)
		-- (0.4cm+0.1cm,-0.1cm-0.4cm)
		-- (-0.4cm-0.1cm,-0.1cm-0.4cm)
		-- cycle;},
	invisible/.style={opacity=0},
	visible on/.style={alt={#1{}{fill=white,white,opacity=1,invisible}}},
	alt/.code args={<#1>#2#3}{%
		\alt<#1>{\pgfkeysalso{#2}}{\pgfkeysalso{#3}} % \pgfkeysalso doesn't change the path
	},
}

\usepackage{fontawesome5}

\newcounter{resnode}

\usepackage{pgfplots}
\usepackage{pgfplotstable}
\pgfplotsset{compat=1.18}
\usepgfplotslibrary{fillbetween,groupplots,dateplot,units,statistics}

\definecolor{green1}{HTML}{008060}
\definecolor{green2}{HTML}{4FC27F}
\definecolor{green3}{HTML}{94BB8E}
\definecolor{green4}{HTML}{B6E6AE}

\definecolor{red1}{HTML}{710000}
\definecolor{red2}{HTML}{b70009}
\definecolor{red3}{HTML}{d35328}
\definecolor{red4}{HTML}{ef8a53}
\definecolor{red5}{HTML}{f9d3bf}

\definecolor{blue1}{HTML}{124066}
\definecolor{blue2}{HTML}{2279BF}
\definecolor{blue3}{HTML}{2EA1FF}
\definecolor{blue4}{HTML}{A7C7FD}

\definecolor{nmse0}{HTML}{006eda}
\definecolor{nmse1}{HTML}{ffffff}

\usepackage{animate}
\usepackage{siunitx}

\sisetup{
	per-mode=symbol,
	%fraction-function=\frac
}

\ExplSyntaxOn
\DeclareExpandableDocumentCommand\myeval{m}{\fp_eval:n{#1}}
\ExplSyntaxOff

\newcommand\myfpnum[1]{\expandafter\num[round-precision=3,round-mode=places]{\myeval{#1}}}

    \pgfplotsset{
        compat=newest,
        /pgfplots/legend image code/.code={%
            \draw[mark repeat=2,mark phase=2,line width=2.5pt,#1]
                plot coordinates {
                    (0cm,0cm)
                    (0.3cm,0cm)
                    (0.6cm,0cm)%
                };
        },
    }


\usepackage{contour}

\newlength{\bucketsep}
\setlength{\bucketsep}{0.75cm}

\tikzstyle{myplotstyle} = [line width=1pt,blue2]

\usepackage{ccicons}

\pgfplotsset{layers/standard/.define layer set={background,
		axis background,axis grid,axis ticks,axis lines,axis tick labels,pre
		main,main,
		axis descriptions,axis foreground
	}{
		grid style={/pgfplots/on layer=axis grid},
		tick style={/pgfplots/on layer=axis ticks},
		axis line style={/pgfplots/on layer=axis lines},
		label style={/pgfplots/on layer=axis descriptions},
		legend style={/pgfplots/on layer=axis descriptions},
		title style={/pgfplots/on layer=axis descriptions},
		colorbar style={/pgfplots/on layer=axis descriptions},
		ticklabel style={/pgfplots/on layer=axis tick labels},
		axis background@ style={/pgfplots/on layer=axis background},
		3d box foreground style={/pgfplots/on layer=axis foreground},
}}

% https://tex.stackexchange.com/questions/55806/mindmap-tikzpicture-in-beamer-reveal-step-by-step
\tikzset{
    invisible/.style={opacity=0},
    visible on/.style={alt={#1{}{invisible}}},
    alt/.code args={<#1>#2#3}{%
      \alt<#1>{\pgfkeysalso{#2}}{\pgfkeysalso{#3}} % \pgfkeysalso doesn't change the path
    },
  }

\pgfplotscreateplotcyclelist{myclist}{
black,every mark/.append style={fill=red1},mark=*\\
black,every mark/.append style={fill=red2},mark=*\\
black,every mark/.append style={fill=red3},mark=*\\
black,every mark/.append style={fill=red4},mark=*\\
black,every mark/.append style={fill=red5},mark=*\\
black,every mark/.append style={fill=blue1},mark=*\\
black,every mark/.append style={fill=blue2},mark=*\\
black,every mark/.append style={fill=blue3},mark=*\\
}

\usepackage{media9}

\newcounter{boxpos}

\setbeamercolor{background canvas}{bg=white}

    \newlength\imwidth
    \setlength{\imwidth}{4cm}

\newcommand{\plotnarma}[4]{
	\addplot[
	fill={#3},
	color={#3},
	%      postaction={
	%        pattern=#4,
	%        pattern color=white,
	%      }
	] table[x=index, y=y_#1, y error=yerr_#1, col sep=comma] {\df}; %
	\addlegendentry{\color{fg}#2}
	\label{nmse_plot_#2}
}

\pgfplotsset{
legend style={fill=BgColor},
}

\pgfdeclarelayer{background}
\pgfdeclarelayer{foreground}
\pgfsetlayers{background,main,foreground}

\usepackage[most,skins]{tcolorbox}
\newtcolorbox{nextstep}[1][]{
    enhanced,
    boxrule=0pt,
    frame hidden,
    sharp corners,
    left skip=-4mm,
    right skip=-2mm,
    left=1mm,
    right=1mm,
    bottom=0mm,
    leftrule=2mm,
    boxsep=3mm,
    borderline west={2mm}{0pt}{SeqBlue3},
    colback=SeqBlue4,
    colbacktitle=SeqBlue4,
    fonttitle=\sffamily\bfseries,
    coltitle=black,
    %title={#2},
    fontupper=\Large
    #1,
}
\newtcolorbox{conclusion}[1][]{
    enhanced,
    boxrule=0pt,
    frame hidden,
    sharp corners,
    left skip=-4mm,
    right skip=-2mm,
    left=1mm,
    right=1mm,
    bottom=0mm,
    leftrule=2mm,
    boxsep=3mm,
    borderline west={2mm}{0pt}{SeqGreen3},
    colback=SeqGreen4,
    colbacktitle=SeqGreen4,
    fonttitle=\sffamily\bfseries,
    coltitle=black,
    %title={#2},
    fontupper=\Large
    #1,
}

% Style to select only points from #1 to #2 (inclusive)
\pgfplotsset{select coords between index/.style 2 args={
    x filter/.code={
        \ifnum\coordindex<#1\def\pgfmathresult{}\fi
        \ifnum\coordindex>#2\def\pgfmathresult{}\fi
    }
}}

\begin{document}
\setbeamercolor{background canvas}{bg=BgColor}

\begin{frame}[plain,label=title]
\thispagestyle{empty}
	\setlength{\imwidth}{2cm}
	\begin{tikzpicture}[overlay,remember picture]
        \matrix[%
            anchor=south west,
            matrix of nodes,
            inner sep=0pt,
            inner sep=0cm,
			column sep=0.2cm,
			row sep=0.2cm,
            ampersand replacement=\&,
        ]  at ($(current page.south west)+(0.25,0.25)$) (m) {
            \includegraphics[width=\imwidth]{figures/ugent_logo} \&
            \includegraphics[width=\imwidth]{figures/ilvo} \&
            \includegraphics[width=\imwidth]{figures/imec_logo} \\
        };

        \matrix[%
            anchor=south east,
            matrix of nodes,
            inner sep=0pt,
            inner sep=0cm,
			column sep=0.2cm,
			row sep=0.2cm,
            ampersand replacement=\&,
        ]  at ($(current page.south east)+(-0.25,0.25)$) (m) {
            \includegraphics[width=\imwidth]{figures/AIRO} \\
        };
	\begin{pgfonlayer}{background}
		\node[opacity=0.3] at ($(current page.center)+(0,0.5)$) {\includegraphics[height=\paperheight]{figures/bg-image}};
	\end{pgfonlayer}
	\end{tikzpicture}
    \maketitle
\end{frame}

\begin{frame}{Plants Are Not Part of Horticulture Environment Control}

\centering
	\begin{tikzpicture}
	
	\node (greenhouse) {\includegraphics[width=0.5\textwidth]{figures/serre}};
	\node[anchor=north east] at ($(greenhouse.north west)+(-0.25,0)$) {\includegraphics[width=0.2\textwidth]{figures/farmer}};
   %\node[anchor=south east] at ($(greenhouse.south west)+(-0.25,0)$) {\Huge\faBrain};
	\node[anchor=north west] at ($(greenhouse.north east)+(0.25,0)$) {\includegraphics[width=0.2\textwidth]{figures/refinery}};
   \node[anchor=south west] at ($(greenhouse.south east)+(0.25,0)$) {\Huge\faCar\hspace{0.5cm}\faTemperatureHigh};
	
	
\end{tikzpicture}
	


\end{frame}

\begin{frame}{Complex Input-driven Responses of Plants}
    \centering

	\setlength{\imwidth}{4.5cm}
    \begin{tikzpicture}%[remember picture, overlay]
        \matrix[%
            anchor=south,
            matrix of nodes,
            ampersand replacement=\&,
            inner sep=0pt,
            inner sep=0cm,
			column sep=0.2cm,
			row sep=0.2cm,
        ] (m) {
            \includegraphics[width=\imwidth]{figures/Mimosa pudica2} \&
            \includegraphics[width=\imwidth]{figures/plant_in_sun.png} \&
            \includegraphics[width=\imwidth]{figures/stomata} \\
            %\includegraphics[width=\imwidth]{figures/bladdikte} \\
        };

		\usebeamercolor{normal text}
		\begin{scope}[
			every node/.style={fill=white,fill opacity=0.7,text opacity=1,text=fg,inner sep=2pt,outer sep=0pt,Text Depth=labels}
		]
		\node[above=0cm of m-1-1.south] {movement};
		\node[above=0cm of m-1-2.south] {colour and day length};
		\node[above=0cm of m-1-3.south] {stomata};
		%\node[above=0cm of m-2-2.south] {bladdikte};
		\end{scope}
    \end{tikzpicture}

    % © Volodymyr Koman/MIT Chemical Engineering
    % https://www.sonnenseite.com/en/science/sensors-applied-to-plant-leaves-warn-of-water-shortage/
\end{frame}

\begin{frame}{Plants Can Be Considered Dynamic Systems}
    \centering
    \setlength{\imwidth}{3cm}
    \begin{tikzpicture}

        \begin{scope}[local bounding box=image]
            \node[anchor=center] (plant) {\includegraphics[width=7cm]{figures/random-plant2}};
        \end{scope}
		\node at (-4.5,0) (input) {};
		\node at (4.5,0) (output) {};

       	\begin{scope}[transparency group,opacity=0.85]
       		\draw[line width=5mm,{Latex[length=10mm, width=12.5mm]}-] (image.center)++(-0.5,0) -- ++(-2,0);
       		\draw[line width=5mm,-{Latex[length=10mm, width=12.5mm]}] (image.center)++(0.5,0) -- ++(2,0);
       	\end{scope}

        \matrix[%
            anchor=center,
            matrix of nodes,
            inner sep=0pt,
            inner sep=0cm,
			column sep=0.2cm,
			row sep=0.2cm,
            ampersand replacement=\&,
        ] at (input) (min) {
            \includegraphics[width=\imwidth]{figures/weather} \\
            \includegraphics[width=\imwidth]{figures/bladluis} \\
            \includegraphics[width=\imwidth]{figures/butterfly} \\
        };

        \matrix[%
            anchor=center,
            matrix of nodes,
            inner sep=0pt,
            inner sep=0cm,
			column sep=0.2cm,
			row sep=0.2cm,
            ampersand replacement=\&,
        ] at (output) (mout) {
            \includegraphics[width=\imwidth]{figures/fall} \\
            \includegraphics[width=\imwidth]{figures/fruit} \\
            \includegraphics[width=\imwidth]{figures/wortels} \\
        };


    \end{tikzpicture}

% image by Marco Verch CC-BY, https://www.flickr.com/photos/30478819@N08/32661392157
\end{frame}



\begin{frame}{Reservoir Computing Enables Computing with Dynamic Systems}
    \centering
    \begin{tikzpicture}
    
    	% draw rectangle labels
    	\node[box label,anchor=base] at (1,3.5) (i-label) {input};
    	\node[box label,anchor=base] at (5,3.5) (res-label) {reservoir/RNN};
    	\node[box label,anchor=base] at (9,3.5) (o-label) {target};
    	
    	\draw[line] (i-label.base) ++(-1,-0.2) -- ++(2,0);
    	\draw[line] (res-label.base) ++(-2.5,-0.2) -- ++(5,0);
    	\draw[line] (o-label.base) ++(-1,-0.2) -- ++(2,0);
    	
    	
    	%%% other reservoir
  	\begin{scope}
      \node[input node] at (1,2) (in0) {};
      \node[input node] at (1,1) (in1) {};

    	\setcounter{resnode}{0}
    	% reservoir nodes
    	\foreach \x/\y in {3/2.25,3.5/1,4/3.5,5/1.5,5.25/3.75,6/0.5,6.15/2.25,7/3.5} {
    		\draw (\x,2.5*\y/3.5) node[reservoir node] (rn\the\value{resnode}) {};
    		\stepcounter{resnode}
    	}

    	% input arrows
    	\foreach \i/\j in {0/0,0/1,1/0,1/2} {
    		\draw[line>] (in\i) -- (rn\j);
    	}

    	% reservoir arrows
    	\foreach \i/\j in {0/3,0/2,1/0,1/5,2/1,3/1,2/3,4/2,5/6,6/2,6/4,6/7,7/0,7/4} {
    		\draw[line>] (rn\i) -- (rn\j);
    	}
    	% bent arrows
        \foreach \i/\j in {2/0,5/7} {
       		\draw[line>] (rn\i) to[bend right] (rn\j);
       	}
       	
    
        	% output nodes
    	\node[output node] at (9,1.5) (out0) {};

    	% output arrows
    	\draw[line>,red2] (rn6) -- (out0);
       	\draw[line>,red2] (rn3) -- (out0);
       \draw[line>,red2] (rn5) -- (out0);
       \draw[line>,red2] (rn7) -- (out0);
       \end{scope}
	\end{tikzpicture}
	
	\tikz{
	\matrix[
		matrix of nodes,
		column sep=0.2cm,
		ampersand replacement=\&,
		column 1/.style={visible on=<2->},
		column 2/.style={visible on=<3->},
		column 3/.style={visible on=<4->},
		column 4/.style={visible on=<5->},
		] {
	\includegraphics[height=2cm]{figures/water_bucket_ripples} \&
	\includegraphics[height=2cm]{figures/nasa-ball-robot} \&
	\includegraphics[height=2cm]{figures/photinic_ic} \&
	\includegraphics[height=2cm]{figures/strawberry}\\
	};
	}
	
\end{frame}


\begin{frame}{Computing With a Water Tank}

\centering
\begin{tikzpicture}%[overlay,remember picture]
\node at (0,0) (bucket) {\includegraphics[width=0.5\linewidth]{figures/water-bucket}};

\node[left=0.5cm of bucket,visible on=<2->] (motor) {\includegraphics[width=1.5cm]{figures/motor}};
\node[right=0.5cm of bucket,visible on=<3->] (camera) {\Huge\faCamera};
\node[left=0.5cm of motor,align=center,font=\Huge,visible on=<2->] (digits) {zero\\one\\one\\zero\\zero\\\faVolumeUp};

\node[right=0.5cm of camera,align=center,font=\Huge,visible on=<3->] (digits2) {0\\1\\1\\0\\0};

\draw[line>,visible on=<2->] (digits) -- (motor);
\draw[line>,visible on=<2->] (motor) -- (bucket);
\draw[line>,visible on=<3->] (bucket) -- (camera);
\draw[line>,visible on=<3->] (camera) -- (digits2);
\end{tikzpicture}

\begin{tikzpicture}[overlay,remember picture]
\node[anchor=west] at ($(current page.south west)+(0.5,0.5)$) {Fernando et al., 2003};
\end{tikzpicture}

\end{frame}


\begin{frame}{Reservoir Computing with Plants}
    \centering
    \begin{tikzpicture}

	\matrix[
        matrix of nodes,
        anchor=center,
        inner sep=0.2em,
		ampersand replacement=\&,
		nodes={anchor=center},
    ] at (-4.5,0) (input) {
        {\color{SeqRed2}\Huge\faTemperatureLow} \& temperature \\
        {\color{SeqBlue2}\Huge\faTint}          \& humidity \\
        {\color{orange}\Huge\faSun}         \& light \\
    };

        \begin{scope}[local bounding box=image,yshift=-2.3cm,xshift=-3.5cm]
            \node[anchor=south west,inner sep=0pt] at (0,0) (plant) {\includegraphics[width=7cm]{figures/strawberry_transp}};

            %\onslide<2->{
            \begin{scope}[x={(plant.south east)},y={(plant.north west)}]
                \draw[fill=red3] (0.7,0.73) circle (3pt) coordinate (m1);
                \draw[fill=red3] (0.3,0.73) circle (3pt) coordinate (m2);
                \draw[fill=red3] (0.325,0.875) circle (3pt) coordinate (m3);
                \draw[fill=red3] (0.7,0.4) circle (3pt) coordinate (m4);
                \draw[fill=red3] (0.39,0.42) circle (3pt) coordinate (m5);
                \draw[fill=red3] (0.75,0.15) circle (3pt) coordinate (m6);
                \draw[fill=red3] (0.41,0.3) circle (3pt) coordinate (m7);
                \draw[fill=red3] (0.55,0.15) circle (3pt) coordinate (m8);
                \draw[fill=red3] (0.18,0.42) circle (3pt) coordinate (m9);
            \end{scope}
            %}
        \end{scope}

	\matrix[
        matrix of nodes,
        anchor=center,
        inner sep=0.2em,
		ampersand replacement=\&,
		nodes={anchor=center},
    ] at (4.5,0) (output) {
        {\color{SeqRed2}\Huge\faTemperatureLow} \& temperature \\
        {\color{SeqBlue2}\Huge\faTint}          \& humidity \\
        {\color{orange}\Huge\faSun}         \& light \\
		{\includegraphics[width=2cm]{figures/IRGA2}} \& \node[align=center]{transpiration\\ photosynthesis}; \\
        {\Huge\faIcon{tachometer-alt}} \& \node[align=center]{NARMA\\ delay line}; \\
    };

        \node[below=0cm of image] (lt) {leaf thickness measurement};
		\draw[fill=red3] (lt.west) circle (3pt);


	\begin{scope}[transparency group,-{Latex[length=10mm, width=12.5mm]},,opacity=0.3]
		\draw[line width=5mm] ($0.5*(input.east)+0.5*(input.center)$) -- ($0.5*(image.west)+0.5*(image.center)$);
		\draw[line width=5mm] ($0.5*(image.east)+0.5*(image.center)$) -- ($0.5*(output.west)+0.5*(output.center)$);
	\end{scope}
    \end{tikzpicture}
    
    \begin{tikzpicture}[overlay,remember picture]
	  \node[fill=bg,text=black,anchor=south west] at ($(current page.south west)+(0.5,0.5)$) (qr) {\qrcode[height=1.75cm]{https://www.nature.com/articles/s41598-022-16874-0}};
	  \node[right=0.5cm of qr.south east,anchor=south west] {Pieters et al., 2022};
    \end{tikzpicture}
\end{frame}


\begin{frame}{Reservoir Computing Applied to FSPMs}
    \centering
    \begin{tikzpicture}
    	% rectangles
    	\draw[box] (0,0) rectangle (2,6);
    	\draw[box] (2.5,0) rectangle (7.5,6);
    	\draw[box] (8,0) rectangle (10,6);

    	% draw rectangle labels
    	\node[box label,anchor=base] at (1,6.2) {input};
    	\node[box label,anchor=base] at (5,6.2) {reservoir};
    	\node[box label,anchor=base] at (9,6.2) {target};
    	
    	%% draw plant reservoir
    	
    	% draw sun, rain and bug
    	\node at (0.5,4) {\color{orange}\huge\faSun};
    	\node at (1.5,4) {\color{blue}\huge\faCloudShowersHeavy};
    	\node at (1,5) {\huge\faBug};
    	

    	% draw plant
    	\draw[postaction={decorate,
    		decoration={
    			markings,
    			mark=between positions 0.2 and 0.9 step 0.45 with {
    				\draw[line width=1pt,SeqGreen2,fill=SeqGreen3,line cap=round] (0,0) to[bend right,looseness=1.2] (.575,-0.46);
    				\draw[line width=1pt,SeqGreen2,fill=SeqGreen3,line cap=round] (0,0) to[bend left,looseness=1.2] (.575,-0.46);
    				\draw[line width=1pt,SeqGreen2,line cap=round] (0,0) -- (0.425,-0.34);
        	}}},
        	postaction={decorate,
        		decoration={
        			markings,
        			mark=between positions 0.4 and 0.5 step 0.2 with {
        				\draw[line width=1pt,SeqGreen2,fill=SeqGreen3,line cap=round] (0,0) to[bend right,looseness=1.2] (.575,0.46);
        				\draw[line width=1pt,SeqGreen2,fill=SeqGreen3,line cap=round] (0,0) to[bend left,looseness=1.2] (.575,0.46);
        				\draw[line width=1pt,SeqGreen2,line cap=round] (0,0) -- (0.425,0.34);
        			}
        	}},
        	SeqGreen2,line width=1.5pt,line cap=round]
    	   (5,4) to[out=90,in=-75] ++(-0.5,2);
    	% draw pot
    	\draw[brown,fill=brown] (5,3.25) -- ++(-0.5,0) -- ++(-0.1,0.75) -- ++(1.2,0) -- ++(-0.1,-0.75) -- cycle;

    	% output    	
    	\draw (7,3.2) node[anchor=center] (leaf) {\Large\color{SeqGreen2}\faLeaf};
        \draw (7,4.1) node[anchor=center] (rh) {\Large\color{SeqRed2}\faApple*};
        \draw (7,4.9) node[anchor=center] (carrot) {\Large\color{SeqYellow2}\faCarrot};
        \draw (7,5.6) node[anchor=center] (carrot) {\Large\color{SeqYellow2}\color{SeqGrey2}\faEyeDropper};
        
       % target
       \only<1>{
       \node[output node] at (9,4.5) (fspm-target) {};
       }
       \only<2>{
    	\node at (9,4) {\color{orange}\huge\faSun};
    	\node at (9,5) {\color{red2}\huge\faTemperatureLow};
    	}
        

    	\draw[line>] (1.9,4.5) -- (4.5,4.5);
    	\draw[line>] (5.5,4.2) -- (6.6,3.2);
    	\draw[line>] (5.5,4.4) -- (6.6,4.1);
    	\draw[line>] (5.5,4.6) -- (6.6,4.9);
    	\draw[line>] (5.5,4.8) -- (6.6,5.6);
    	%\only<1>{\draw[line>] (7.3,4.5) -- (8.5,4.5) node[midway,above,text=fg,fill=BgColor,fill opacity=0.90,text opacity=1] {\Large ?};}
    	\only<1>{\draw[line>] (7.3,4.5) -- (8.5,4.5);}
    	%\draw[line>] (7,-2.25) -- (rh);
    	
    	%%% other reservoir
  	\begin{scope}
      \node[input node] at (1,2) (in0) {};
      \node[input node] at (1,1) (in1) {};

    	\setcounter{resnode}{0}
    	% reservoir nodes
    	\foreach \x/\y in {3/2.25,3.5/1,4/3.5,5/1.5,5.25/3.75,6/0.5,6.15/2.25,7/3.5} {
    		\draw (\x,2.5*\y/3.5) node[reservoir node] (rn\the\value{resnode}) {};
    		\stepcounter{resnode}
    	}

    	% input arrows
    	\foreach \i/\j in {0/0,0/1,1/0,1/2} {
    		\draw[line>] (in\i) -- (rn\j);
    	}

    	% reservoir arrows
    	\foreach \i/\j in {0/3,0/2,1/0,1/5,2/1,3/1,2/3,4/2,5/6,6/2,6/4,6/7,7/0,7/4} {
    		\draw[line>] (rn\i) -- (rn\j);
    	}
    	% bent arrows
        \foreach \i/\j in {2/0,5/7} {
       		\draw[line>] (rn\i) to[bend right] (rn\j);
       	}
       	
    
        	% output nodes
    	\node[output node] at (9,1.5) (out0) {};

    	% output arrows
    	\draw[line>,red2] (rn6) -- (out0);
       	\draw[line>,red2] (rn3) -- (out0);
       \draw[line>,red2] (rn5) -- (out0);
       \draw[line>,red2] (rn7) -- (out0);
       \end{scope}
	\end{tikzpicture}
\end{frame}

\begin{frame}{Selecting FSPM Reservoir Readout and Target}
    \centering
    \begin{tikzpicture}
    	% rectangles
    	\draw[box] (0,0) rectangle (2,6);
    	\draw[box] (2.5,0) rectangle (7.5,6);
    	\draw[box] (8,0) rectangle (10,6);

    	% draw rectangle labels
    	\node[box label,anchor=base] at (1,6.2) {input};
    	\node[box label,anchor=base] at (5,6.2) {reservoir};
    	\node[box label,anchor=base] at (9,6.2) {target};
    	
    	%% draw plant reservoir
    	
    	% draw sun, rain and bug
    	\node at (0.5,4) {\color{orange}\huge\faSun};
    	\node at (1.5,4) {\color{blue}\huge\faCloudShowersHeavy};
    	\node at (1,5) {\huge\faBug};
    	

    	% draw plant
    	\draw[postaction={decorate,
    		decoration={
    			markings,
    			mark=between positions 0.2 and 0.9 step 0.45 with {
    				\draw[line width=1pt,SeqGreen2,fill=SeqGreen3,line cap=round] (0,0) to[bend right,looseness=1.2] (.575,-0.46);
    				\draw[line width=1pt,SeqGreen2,fill=SeqGreen3,line cap=round] (0,0) to[bend left,looseness=1.2] (.575,-0.46);
    				\draw[line width=1pt,SeqGreen2,line cap=round] (0,0) -- (0.425,-0.34);
        	}}},
        	postaction={decorate,
        		decoration={
        			markings,
        			mark=between positions 0.4 and 0.5 step 0.2 with {
        				\draw[line width=1pt,SeqGreen2,fill=SeqGreen3,line cap=round] (0,0) to[bend right,looseness=1.2] (.575,0.46);
        				\draw[line width=1pt,SeqGreen2,fill=SeqGreen3,line cap=round] (0,0) to[bend left,looseness=1.2] (.575,0.46);
        				\draw[line width=1pt,SeqGreen2,line cap=round] (0,0) -- (0.425,0.34);
        			}
        	}},
        	SeqGreen2,line width=1.5pt,line cap=round]
    	   (5,4) to[out=90,in=-75] ++(-0.5,2);
    	% draw pot
    	\draw[brown,fill=brown] (5,3.25) -- ++(-0.5,0) -- ++(-0.1,0.75) -- ++(1.2,0) -- ++(-0.1,-0.75) -- cycle;

    	% output    	
    	\draw (7,3.2) node[anchor=center] (leaf) {\Large\color{SeqGreen2}\faLeaf};
        \draw (7,4.1) node[anchor=center] (rh) {\Large\color{SeqGreen1}\faLeaf};
        \draw (7,4.9) node[anchor=center] (carrot) {\Large\color{SeqGreen3}\faLeaf};
        \draw (7,5.6) node[anchor=center] (carrot) {\Large\color{SeqGreen4}\faLeaf};
        
       % target
    	\node at (9,4) {\color{orange}\huge\faSun};
    	\node at (9,5) {\color{red2}\huge\faTemperatureLow};
        

    	\draw[line>] (1.9,4.5) -- (4.5,4.5);
    	\draw[line>] (5.5,4.2) -- (6.6,3.2);
    	\draw[line>] (5.5,4.4) -- (6.6,4.1);
    	\draw[line>] (5.5,4.6) -- (6.6,4.9);
    	\draw[line>] (5.5,4.8) -- (6.6,5.6);
    	
    	\draw[line>] (7.3,4.5) -- (8.5,4.5) node[midway] (h1) {};
    	
    	\draw[fill=red2,line,red2] (h1) circle (0.1cm);
      \draw[red2,line] (h1) |- ++(-1,-2.5) node[anchor=east,overlap label] (rr) {\color{red2} Ridge regression};
    	
    	\node[below right=0.5cm and 0.5cm of rr.south,overlap label] (rrlr) {linear regression};
    	\node[below left=0.5cm and 0.5cm of rr.south,overlap label] (rrpo) {prevents overfit};
    	
    	\draw[line>] (rr) -- (rrlr);
    	\draw[line>] (rr) -- (rrpo);
    \end{tikzpicture}
\end{frame}

\begin{frame}{Overview of the Methodology (Practical Point of View)}

\centering

\begin{columns}
\begin{column}{0.47\linewidth}
\centering
\begin{tikzpicture}[every node/.style={draw,rounded corners,align=center}]
\node[] (run) {run the FSPM};
\node[below=0.5cm of run] (data) {collect output data\\into table by unit};
\node[below=0.5cm of data] (prep) {data preprocessing};
\node[below=0.5cm of prep] (train) {train Ridge regression};
\node[below=0.5cm of train] (eval) {evaluate on unseen data};

\draw[line>] (run) -- (data);
\draw[line>] (data) -- (prep);
\draw[line>] (prep) -- (train);
\draw[line>] (train) -- (eval);
\draw[line>] (prep.west) -- ++(-1,0) |- (eval);
\end{tikzpicture}
\end{column}%
\hfill%
\begin{column}{0.47\linewidth}
\centering
\only<2->{
\begin{tikzpicture}
    \begin{axis}[
    	width=\linewidth,
    	height=\linewidth,
    	date coordinates in=x,
    	xticklabel={\month-\day},
    	xlabel={time [month-day]},
    	ylabel={air temperature [\si{\degreeCelsius}]},
    	xtick={2012-06-01,2012-06-05,2012-06-09,2012-06-13},
    	]
        \addplot[line width=1pt] table[x=time, y=Tac, col sep=semicolon] {data/meteo.csv};
        
        \scoped[on background layer] \fill [blue3!50] (axis cs:2012-06-01 00:00,\noexpand\pgfkeysvalueof{/pgfplots/ymin}) rectangle (axis cs:2012-06-05 00:00,\noexpand\pgfkeysvalueof{/pgfplots/ymax});
        
        \scoped[on background layer] \fill [green2!50] (axis cs:2012-06-05 00:00,\noexpand\pgfkeysvalueof{/pgfplots/ymin}) rectangle (axis cs:2012-06-09 00:00,\noexpand\pgfkeysvalueof{/pgfplots/ymax});
        
        \scoped[on background layer] \fill [blue3!50] (axis cs:2012-06-09 00:00,\noexpand\pgfkeysvalueof{/pgfplots/ymin}) rectangle (axis cs:2012-06-13 00:00,\noexpand\pgfkeysvalueof{/pgfplots/ymax});
        
        \scoped[on background layer] \fill [green2!50] (axis cs:2012-06-13 00:00,\noexpand\pgfkeysvalueof{/pgfplots/ymin}) rectangle (axis cs:2012-06-14 00:00,\noexpand\pgfkeysvalueof{/pgfplots/ymax});
        
        \draw (axis cs: 2012-06-03 00:00,13) node {train};
        \draw (axis cs: 2012-06-07 00:00,13) node {test};
        \draw (axis cs: 2012-06-11 00:00,13) node {train};
    \end{axis}
\end{tikzpicture}
}
\end{column}
\end{columns}
\end{frame}

\begin{frame}{Normalised Evaluation Metric}

$ \text{NMSE} = \dfrac{1}{N}\dfrac{\sum_{i=0}^{N-1}\left[y-\hat{y}\right]^2}{\text{var}(y)} $

\begin{tabular}{cl}
+ & normalised metric \\
+ & easy to compare different targets \\
- & phase/time sensitive \\
- & sensitive to outliers
\end{tabular}

\end{frame}

\begin{frame}{Simple Experimental Evaluation Using Three FSPMs}

\centering
\begin{tabular}{ccc}
\includegraphics[width=0.3\textwidth]{figures/cnwheat} & \includegraphics[width=0.3\textwidth]{figures/hydroshoot} & \includegraphics[width=0.3\textwidth]{figures/GrassLeaf} \\
\textbf{WheatFspm} & \textbf{Hydroshoot} & \textbf{GrassLeaf} \\
{\footnotesize (Barillot et al., 2016)} & {\footnotesize (Albasha et al., 2019)} & {\footnotesize (unpublished)} \\
transpiration & transpiration & leaf thickness \\
\end{tabular}

\end{frame}
\begin{frame}{Transpiration-based Reservoirs Based on Hydroshoot}
\centering
\begin{tikzpicture}
\begin{axis}[
		name=graph1,
    	%width=\linewidth,
    	%height=\linewidth,
    	xlabel={time [day]},
    	ylabel={air temperature [\si{\degreeCelsius}]},
    	legend pos=outer north east,
    	legend style={cells={align=left}},
    	]
    	
    \addplot[line width=1pt,red2,select coords between index={1}{19}] table[x=time, y=y, col sep=comma] {data/results_Hydroshoot-Tair.csv};  	
    \addlegendentry{\color{fg} $y$};
  \foreach \x/\y in {20/39,40/59,60/79,80/99,100/119} {
  \addplot[line width=1pt,red2,select coords between index={\x}{\y},forget plot] table[x=time, y=y, col sep=comma] {data/results_Hydroshoot-Tair.csv};}
    	
  \addplot[line width=1pt,blue3,select coords between index={1}{19}] table[x=time, y=yh, col sep=comma] {data/results_Hydroshoot-Tair.csv};
  \addlegendentry{\color{fg} $\hat{y}$};
    	
\foreach \x/\y in {20/39,40/59,60/79,80/99,100/119} {
  \addplot[line width=1pt,blue3,select coords between index={\x}{\y},forget plot] table[x=time, y=yh, col sep=comma] {data/results_Hydroshoot-Tair.csv};
  }
  

  
  

\end{axis}
\draw (graph1.south east) node[anchor=south west] {NMSE=0.62};
\end{tikzpicture}
\end{frame}

\begin{frame}{Transpiration-based Reservoirs Based on WheatFspm}
\centering
\begin{tikzpicture}
\begin{axis}[
		name=graph1
    	%width=\linewidth,
    	%height=\linewidth,
    	xlabel={time [day]},
    	ylabel={air temperature [\si{\degreeCelsius}]},
    	legend pos=outer north east,
    	legend style={cells={align=left}},
    	]
    	
  \addplot[line width=1pt,red2,select coords between index={1}{19}] table[x=time, y=y, col sep=comma] {data/results_WheatFspm-Tair.csv};
  \addlegendentry{\color{fg} $y$};
    	\foreach \x/\y in {20/39,40/59,60/79,80/99,100/119} {
  \addplot[line width=1pt,red2,select coords between index={\x}{\y},forget plot] table[x=time, y=y, col sep=comma] {data/results_WheatFspm-Tair.csv};
  }
  
    \addplot[line width=1pt,blue2] table[x=time, y=yh, col sep=comma,select coords between index={1}{19}] {data/results_WheatFspm-Tair.csv};
    \addlegendentry{\color{fg} $\hat{y}$};
      	\foreach \x/\y in {20/39,40/59,60/79,80/99,100/119} {
  \addplot[line width=1pt,blue2] table[x=time, y=yh, col sep=comma,select coords between index={\x}{\y},forget plot] {data/results_WheatFspm-Tair.csv};
  }
  
\end{axis}
\draw (graph1.south east) node[anchor=south west] {NMSE=0.19};
\end{tikzpicture}
\end{frame}

\begin{frame}{Transpiration-based Reservoirs Based on GrassLeaf}
\centering
\begin{tikzpicture}
\begin{axis}[
name=graph1
    	%width=\linewidth,
    	%height=\linewidth,
    	xlabel={time [day]},
    	ylabel={air temperature [\si{\degreeCelsius}]},
    	legend pos=outer north east,
    	legend style={cells={align=left}},
    	]
  \addplot[line width=1pt,red2,select coords between index={1}{19}] table[x=time, y=y, col sep=comma] {data/results_GrassLeaf-Tair.csv};
  \addlegendentry{\color{fg} $y$};
  
  \foreach \x/\y in {20/39,40/59,60/79,80/99,100/119,120/139,140/159,160/179} {
  \addplot[line width=1pt,red2] table[x=time, y=y, col sep=comma,select coords between index={\x}{\y},forget plot] {data/results_GrassLeaf-Tair.csv};
  }
  
  \addplot[line width=1pt,blue2,select coords between index={1}{19}] table[x=time, y=yh, col sep=comma] {data/results_GrassLeaf-Tair.csv};
  \addlegendentry{\color{fg} $\hat{y}$};
  \foreach \x/\y in {20/39,40/59,60/79,80/99,100/119,120/139,140/159,160/179} {
  \addplot[line width=1pt,blue2] table[x=time, y=yh, col sep=comma,select coords between index={\x}{\y},forget plot] {data/results_GrassLeaf-Tair.csv};
  }
\end{axis}
\draw (graph1.south east) node[anchor=south west] {NMSE=0.63};
\end{tikzpicture}
\end{frame}

\begin{frame}{Reservoir Readout Error Estimation}
	\centering
    \begin{tikzpicture}
    	\pgfplotstableread[col sep=comma]{data/error_Hydroshoot-PAR.csv}\dfa
    	\pgfplotstableread[col sep=comma]{data/error_WheatFspm-PAR.csv}\dfb
    	\pgfplotstableread[col sep=comma]{data/error_GrassLeaf-PAR.csv}\dfc

    	\begin{axis}[%
    		boxplot,
    		boxplot/draw direction=y,
    		width=0.8\linewidth,
    		height=7cm,
    		typeset ticklabels with strut,
    		ymin=0,
    		ymax=1.0,
    		enlarge x limits={abs=0.5},
    		minor y tick num=1,
    		xtick={1,2,3},
         xticklabels={Hydroshoot,WheatFspm,GrassLeaf},
    		ymajorgrids=true,
    		yminorgrids=true,
    		ylabel={NMSE},
            %legend entries={Hydroshoot,WheatFSPM,GrassLeaf},
            %legend style={at={(0.02,0.98)},anchor=north west},
            xtick style={draw=none},
    	]
        \setcounter{boxpos}{0}

        %\addlegendimage{mark=square*,red1}
        %\addlegendimage{mark=square*,red2}
        %\addlegendimage{mark=square*,blue2}
        
        \addplot[boxplot/draw position=1,red1,line width=1pt,mark=*,] table[y=nmse_test]{\dfa};
        \addplot[boxplot/draw position=2,red1,line width=1pt,mark=*,] table[y=nmse_test]{\dfb};
        \addplot[boxplot/draw position=3,red1,line width=1pt,mark=*,] table[y=nmse_test]{\dfc};
        


    	\end{axis}
    \end{tikzpicture}
\end{frame}

\begin{frame}{Next Steps For This Research}
\begin{nextstep}{}
Create systematic approach to analyse stationarity of readouts
\end{nextstep}
\begin{nextstep}{}
Remove the requirement for stationarity, which is a limiting factor for practical usage
\end{nextstep}
\begin{nextstep}{}
Use FSPMs to Speed-up Greenhouse Experiments
\end{nextstep}
\end{frame}

\begin{frame}{What You Should Remember From This Talk}

\begin{conclusion}{}
FSPMs are non-linear dynamic systems that fit the requirements for reservoir computing
\end{conclusion}

\begin{conclusion}{}
Reservoir computing offers a novel holistic view on plant responses
\end{conclusion}

\end{frame}

\begin{frame}[plain,label=title]
\thispagestyle{empty}
	\setlength{\imwidth}{2cm}
	\begin{tikzpicture}[overlay,remember picture]
        \matrix[%
            anchor=south west,
            matrix of nodes,
            inner sep=0pt,
            inner sep=0cm,
			column sep=0.2cm,
			row sep=0.2cm,
            ampersand replacement=\&,
            nodes={anchor=south},
        ]  at ($(current page.south west)+(0.25,0.25)$) (m) {
            \includegraphics[width=\imwidth]{figures/ugent_logo} \&
            \includegraphics[width=\imwidth]{figures/ilvo} \&
            \includegraphics[width=\imwidth]{figures/imec_logo} \&[1cm]
				\node[fill=white,text=black] {\qrcode[height=\imwidth]{https://github.com/opieters/fspm2023}}; \&[1cm]
				\node[fill=white,text=black] {\qrcode[height=\imwidth]{https://www.nature.com/articles/s41598-022-16874-0}}; \\
        };

        \matrix[%
            anchor=south east,
            matrix of nodes,
            inner sep=0pt,
            inner sep=0cm,
			column sep=0.2cm,
			row sep=0.2cm,
            ampersand replacement=\&,
        ]  at ($(current page.south east)+(-0.25,0.25)$) (m) {
            \includegraphics[width=\imwidth]{figures/AIRO} \\
        };
        
        \matrix[%
            anchor=north,
            matrix of nodes,
            inner sep=0pt,
            inner sep=0cm,
			column sep=0.2cm,
			row sep=0.2cm,
            ampersand replacement=\&,
        ]  at ($(current page.center)+(0,0)$) (m) {
            \includegraphics[width=1cm]{headshots/mcannoodt} \&
            \includegraphics[width=1cm]{headshots/tdeswaef} \&
            \includegraphics[width=1cm]{headshots/mstock} \&
            \includegraphics[width=1cm]{headshots/fwyffels} \\
            Maxime Cannoodt \& Tom De Swaef \& Michiel Stock \& Francis wyffels\\
        };
	\begin{pgfonlayer}{background}
		\node[opacity=0.3] at ($(current page.center)+(0,0.5)$) {\includegraphics[height=\paperheight]{figures/bg-image}};
	\end{pgfonlayer}
	\end{tikzpicture}
	\author{}\date{}\institute{}
    \maketitle
\end{frame}

\end{document}
